\begin{table}[htbp]
\centering
\caption{Robustness: Placebo, Non-Movers, and Urban/Rural Heterogeneity}
\label{tab:robust}
\begin{adjustbox}{max width=\textwidth}
\begin{tabular}{lccccccc}
\hline\hline
 & \multicolumn{3}{c}{Placebo: 1910--1920} & \multicolumn{2}{c}{Non-Movers} & \multicolumn{2}{c}{Urban/Rural} \\
\cmidrule(lr){2-4} \cmidrule(lr){5-6} \cmidrule(lr){7-8}
 & All & Married & Unmarr. & All & Married & Urban M. & Rural M. \\
 & (1) & (2) & (3) & (4) & (5) & (6) & (7) \\
\hline
Exposure (std.) & -0.0010** & 0.0023*** & -0.0078*** & -0.0110*** & -0.0071*** & -0.0062*** & -0.0052*** \\
 & (0.0005) & (0.0006) & (0.0026) & (0.0010) & (0.0008) & (0.0008) & (0.0014) \\
\hline
Controls & Yes & Yes & Yes & Yes & Yes & Yes & Yes \\
State FE & Yes & Yes & Yes & Yes & Yes & Yes & Yes \\
\hline\hline
\end{tabular}
\end{adjustbox}
\begin{tablenotes}[flushleft]
\small
\item \textit{Notes:} Columns 1--3: placebo test using the 1910--1920 linked panel (pre-Johnson-Reed Act). If exposure predicts LFP changes before the Act, the identification strategy fails. Columns 4--5: restrict to women who remained in the same county between 1920 and 1930 (mover $= 0$). Columns 6--7: married women split by urban (non-farm) vs.~rural (farm) residence in 1920. All specifications include individual controls and state fixed effects. Standard errors clustered at county level. $^{***}p<0.01$, $^{**}p<0.05$, $^{*}p<0.1$.
\end{tablenotes}
\end{table}
